\chapter*{PRESUPUESTO}
\renewcommand{\thetable}{P.\arabic{table}}
\setcounter{table}{0}
\label{ch:presupuesto}

\addcontentsline{toc}{chapter}{Presupuesto} % si queremos que aparezca en el índice
\markboth{PRESUPUESTO}{PRESUPUESTO} % encabezado

% TODO (Fase 5): rellenar con las horas reales del diario de desarrollo.
% Fuentes de los precios (coste/hora de referencia, amortización del equipo): citarlas.

En esta sección se presentan los costes asociados al desarrollo del proyecto. Al tratarse de un sistema software, el coste principal corresponde a las horas de trabajo de ingeniería; el resto lo constituyen la amortización del equipo y las licencias, que en este caso son de software libre.

\definecolor{LightCyan}{rgb}{0.88,1,1}

\begin{table}[ht]
\centering
\begin{tabular}{ |m{6.5cm}|c|r|r| }
 \hline
 \rowcolor{LightCyan}
 Concepto & Cantidad & Precio (\euro) & Importe (\euro)\\
 \hline
 Horas de ingeniería (análisis, desarrollo, pruebas, documentación) & -- h & -- & -- \\
 \hline
 Amortización del equipo informático & -- meses & -- & -- \\
 \hline
 Licencias de software (Python, Dash, \LaTeX{}: software libre) & -- & 0,00 & 0,00 \\
 \hline
 \textbf{Total} &  &  & \textbf{--} \\
 \hline
\end{tabular}
\caption{Presupuesto del proyecto.}
\label{tabla:presupuesto}
\end{table}
